\uuid{qFW8}
\titre{Variance d'une transformation affine}
\auteur{Maxime Nguyen}
\organisation{AMSCC}
\datecreate{2026-11-24}
\niveau{L2}
\module{Probabilité}
\chapitre{Probabilité continue}
\sousChapitre{Espérance et Fonction de répartition}
\difficulte{1}
\theme{variance, transformation affine}

\contenu{
\begin{question}
\qtype{onlyone}
On note $\mathbb{V}(X)$ la variance d'une variable $X$. Soit $X$ une variable telle que $\mathbb{V}(X) = 4$. Alors $\mathbb{V}(-2X + 3)$ vaut :
\begin{answers}
\bad{$8$}
\bad{$11$}
\good{$16$}
\bad{$19$}
\end{answers}
\begin{explanations}
Pour tous réels $a$ et $b$, $\mathbb{V}(aX + b) = a^2\,\mathbb{V}(X)$ : la constante
ne change pas la dispersion et le facteur passe au carré.

A : $|-2| \times 4$, le carré a été oublié.

B : $2 \times 4 + 3$, la variance traitée comme une espérance (en corrigeant le signe).

D : $(-2)^2 \times 4 + 3$, la constante a été gardée.

Contre-sens visé : traiter la variance comme l'espérance (linéarité), en
tout ou en partie.
\end{explanations}
\end{question}
}
